\documentclass[11pt,a4paper]{article}
\usepackage[T1]{fontenc}
\usepackage[utf8]{inputenc}
\usepackage{lmodern}
\usepackage[margin=25mm]{geometry}
\usepackage{microtype}
\usepackage{array,booktabs,longtable,needspace}
\usepackage[hidelinks]{hyperref}
\setlength{\parindent}{0pt}
\setlength{\parskip}{0.65em}
\setlength{\tabcolsep}{4pt}
\setlength{\LTleft}{0pt}
\setlength{\LTright}{\fill}
\renewcommand{\arraystretch}{1.25}


\hypersetup{pdftitle={S06 -- Legal Counsel \& Compliance Officers}}

\begin{document}
% S06: extracted from S06 — Legal Counsel and Compliance Officers.tex

{\LARGE\bfseries S06 -- Legal Counsel \& Compliance Officers\par}\vspace{1em}
\noindent\rule{\linewidth}{0.4pt}

\section*{Description}

Legal Counsel and Compliance Officers are individuals or teams responsible for advising the company on applicable legal, regulatory and procedural requirements and overseeing adherence to established obligations. Legal Counsel interprets applicable requirements and advises on legal risks, while Compliance Officers monitor adherence to established requirements. They are grouped together because their interests in GOSPL are closely aligned.

In GOSPL, their focus is the compliance of proposals, voting conditions and actual voting activities. GOSPL supports human compliance assessments and applies established voting conditions; it does not independently interpret legislation or guarantee legal validity.

\section*{Stakeholder Concerns}

\subsection*{Distinct Concerns}

\begingroup\small
\begin{longtable}{@{}>{\raggedright\arraybackslash}p{\dimexpr 0.15\textwidth-2\tabcolsep\relax}>{\raggedright\arraybackslash}p{\dimexpr 0.27\textwidth-2\tabcolsep\relax}>{\raggedright\arraybackslash}p{\dimexpr 0.58\textwidth-2\tabcolsep\relax}@{}}
\toprule
\textbf{ID} & \textbf{Concern} & \textbf{Stakeholder Expectation} \\
\midrule
\endfirsthead
\toprule
\textbf{ID} & \textbf{Concern} & \textbf{Stakeholder Expectation} \\
\midrule
\endhead
\bottomrule
\endfoot
\texttt{S06-C01} & \textbf{Proposal Compliance} & We want to assess whether proposed matters and their supporting documentation are consistent with applicable legal, regulatory and corporate governance requirements. \\
\addlinespace
\texttt{S06-C02} & \textbf{Voting Rule Compliance} & We want to ensure that the conditions governing shareholder eligibility, voting rights, representative authorizations and voting activities are consistent with applicable legal and procedural requirements. \\
\addlinespace
\texttt{S06-C03} & \textbf{Procedural Compliance Assessment} & We want to assess whether voting activities have been conducted according to established procedures, using sufficient evidence to identify deviations, document findings and track the resolution of compliance issues. \\
\end{longtable}
\endgroup

\Needspace{14\baselineskip}

\section*{Potential Stakeholder Conflicts}

\begingroup\small
\begin{longtable}{@{}>{\raggedright\arraybackslash}p{\dimexpr 0.2\textwidth-2\tabcolsep\relax}>{\raggedright\arraybackslash}p{\dimexpr 0.2\textwidth-2\tabcolsep\relax}>{\raggedright\arraybackslash}p{\dimexpr 0.2\textwidth-2\tabcolsep\relax}>{\raggedright\arraybackslash}p{\dimexpr 0.2\textwidth-2\tabcolsep\relax}>{\raggedright\arraybackslash}p{\dimexpr 0.2\textwidth-2\tabcolsep\relax}@{}}
\toprule
\textbf{Conflict} & \textbf{Conflicting Stakeholders} & \textbf{Stakeholder Expectation} & \textbf{Other Stakeholder Expectation} & \textbf{Implication for GOSPL} \\
\midrule
\endfirsthead
\toprule
\textbf{Conflict} & \textbf{Conflicting Stakeholders} & \textbf{Stakeholder Expectation} & \textbf{Other Stakeholder Expectation} & \textbf{Implication for GOSPL} \\
\midrule
\endhead
\bottomrule
\endfoot
\textbf{Proposal Intent vs. Legal Compliance} & S04 -- Board of Directors & Ensure proposals comply with applicable requirements. & Preserve intended proposal objectives. & Support review feedback and authorized amendments. \\
\addlinespace
\textbf{Operational Efficiency vs. Procedural Compliance} & S05 -- AGM Administrators & Ensure applicable procedures are followed. & Execute voting efficiently with minimal administrative overhead. & Support necessary reviews and controls without unnecessary complexity. \\
\addlinespace
\textbf{Voting Flexibility vs. Procedural Restrictions} & S01, S02, S03 -- Shareholders \& Proxies & Keep voting activities within permitted conditions. & Exercise, revise and delegate voting rights with sufficient flexibility. & Communicate and apply established restrictions. \\
\addlinespace
\textbf{Compliance Evidence vs. Privacy} & S01, S02, S03 -- Shareholders \& Proxies & Access and retain sufficient evidence for compliance assessment. & Protect personal information and voting choices from unnecessary access and retention. & Support controlled access, data minimization and justified retention. \\
\end{longtable}
\endgroup

The first two conflicts are reciprocal to previously identified S04 and S05 conflicts. The latter two are retained as proposed conflicts pending explicit confirmation. Compliance Assurance vs. Verification Independence is deferred to S07; it may represent separation of responsibilities rather than a conflict.

\Needspace{10\baselineskip}

\section*{Stakeholder Prioritization}

\begingroup\small
\begin{longtable}{@{}>{\raggedright\arraybackslash}p{\dimexpr 0.23\textwidth-2\tabcolsep\relax}>{\raggedright\arraybackslash}p{\dimexpr 0.77\textwidth-2\tabcolsep\relax}@{}}
\toprule
\textbf{Attribute} & \textbf{Assessment} \\
\midrule
\endfirsthead
\toprule
\textbf{Attribute} & \textbf{Assessment} \\
\midrule
\endhead
\bottomrule
\endfoot
\textbf{Category} & Corporate Governance Stakeholder \\
\addlinespace
\textbf{Justification} & Legal Counsel and Compliance Officers influence the applicable constraints, procedural conditions and compliance expectations governing shareholder voting. Their concerns affect the permissible functional scope and non-functional requirements of GOSPL. \\
\addlinespace
\textbf{Importance} & 5 \\
\addlinespace
\textbf{Influence} & 5 \\
\end{longtable}
\endgroup

\end{document}
